\chapter{Fejlesztői dokumentáció}
\label{ch:impl}

\section{Specifikáció}

A feladat egy webalkalmazás fejlesztése ami személyre szabott önéletrajzokat készít a felhasználónak. A program lehetőséget nyújt a felhasználónak adatai beviltelére szöveges mezőkön, file feltöltésen, illetve LinkedIn profil megadásán keresztül. A program generatív mesterséges intelligencia és Retrieval-Augmented Generation segítségével feldolgozza a felhasználó által megagott adatokat és fileokat, majd egy személyre szabott önéletrajzot készít. A kész önéletrajzot a felhasználó áttekintheti, szükség esetén újrageneráltathatja, illetve több formátumban letöltheti. A felhasználó akármikor visszatérhet az előző lépéshez, vagy megszakíthatja a munkamenetet. Minden adat amit a felhasználó megadott, törlésre kerül a programból való kilépés után, valamint az aktuális munkamenet megszakítását követően.

\subsection{Követelményelemzés}

\subsubsection{Funkcionális követelmények}

\paragraph{Megnyitás:}

A program webalkalmazás formájában elérhető modern webböngészőből (Pl: Google Chrome, Microsoft Edge, Mozzilla Firefox, stb...). Az alkalmazás elérése után a nyitólap(link to nyitólap image) látható.

\paragraph{Adatbevitel:}

A felhasználónak több módon áll lehetősége adatokat bevinni.

\begin{enumerate}

    \item File feltöltéssel
    
A feltöltött fileok lehetnek szöveges vagy kép formátumban, és tartalmazhatnak korábbi önéletrajzokat, okleveleket, elismervényeket, bizonyítványokat és más tanúsítványokat.
    
    \item Szöveges mezőn keresztül

A felhasználó manuálisan megadhat adatokat tanulmányairól, képzéseiről és munkatapasztalatairól. A bevitt szöveges adatok tartalmáról a felhasználónak nyilatkoznia kell lsitából való kiválasztással. 
    
    \item LinkedIn profil link megadásával

A felhasználó Linkedin profiljának linkje segítségével a felhasználó profilja és adatai gyorsan és egyszerűen kitölthetőek.

\end{enumerate} 

\paragraph{Állás specifikálása:}

Az adatbevitel után a felhasználónak specifikálnia kell az állást szöveges bevitelen keresztül. A felhasználó megadhatja az álláshirdetést, vagy egy linket ami a kiíráshoz mutat.

\paragraph{Visszalépés:}

A felhasználónak lehetősége van minden lépés után visszalépni az előző oldalra.

\paragraph{Generálás:}

Az alkalmazás a felhasználó által megadott adatok alapján generál egy önéletrajzot, majd ezt prezentálja a felhasználónak. A felhasználó megismételheti a generálási folyamatot, amennyiben nem elégedett a kimenettel.

\paragraph{Letöltés:}

A generálás után a felhasználó kiválaszthatja, milyen formátumban szeretné letölteni az önéletrajzát. Az elérhető formátumok: PDF, docx, md

\paragraph{Kilépés}

Miután a felhasználó befejezte önéletrajzának elkészítését, kiléphet az alkalmazásból a böngészőablak bezárásával, vagy visszaléphet a kezdőképernyőre az aktuális munkamenet megszakítását követően.

\subsubsection{Nem funkcionális követelmények}

\paragraph{Felhasználók:}

A felhasználónak nem kell regisztrálnia az alkalmazás használatához. Az alkalmazásnak nincsen adminisztrátori felülete.

\paragraph{Munkamenet:}

Az alkalmazás webböngészővel történő megnyitása során az alkalmazás egyedi azonosítóval látja el az aktuális munkamenetet, amit a webböngésző ideiglenesen tárol. Ezen azonosító segítségével követi az alkalmazás a felhasználó által megadott adatokat és feltöltött fileokat. A munkamenet megszakadása történhet automatikusan, 10 perc inaktivitás után, illetve manuálisan, a felhasználó kérésére.

\paragraph{Adatkezelés:}

A felhasználó által megadott adatok és feltöltött dokumentumok nem kerülnek perzisztens tárolásra. Az alkalmazás elindításakor az adatbázis üres állapotba inicializálódik. Az adatbázisban csak az aktív munkamenetekhez köthető adatok kerülnek tárolásra. A munkamenet megszakítása után ezek az adatok törlődnek, és nem lehet őket visszaállítani.

\paragraph{TODO:}

\begin{itemize}
    \item MI-vel kapcsolatos dolgok
    \begin{itemize}
        \item Hallucináció elkerülése
        \item Adatbiztonság
        \item Követelmények beépítése
        \item Promptok specifikálása
        \item Ügynökhálózatok
        \item beépített optimalizációk
    \end{itemize}
    \item Önéletrajz kritériumok
    \item Formai kritériumok
    \item Kimeneti formátum
\end{itemize}

\subsection{Megszorítások}

\begin{itemize} 
    \item Az alkalmazásnak modern webböngészőkőn keresztül, telepítés nélkül kell működnie.
    \item Az üzleti logikát Spring Boot keretrendszerben, Java nyelven kell implementálni
    \item   Az alkalmazás grafikus felületét React keretrendszerben kell elkészíteni
    \item A generatív mesterséges inelligenciát LangFlow keretrendszerben készített ügynökhálózaton keresztül kell elérni
\end{itemize}

\subsection{Használati esetek}

Felhasználói történetek

\section{Tervezés}

\subsection {Diagrammok}

\section{Megvalósítás}

\section{Tesztelés}
